\section{A selection process with unbounded extinction}
\label{sec:selection}

Let $R$ be a unital $k$-algebra, let $e\in R$ be a central idempotent,
and let $\alpha:R\to R$ be a unital endomorphism.  For a left
$R$-module $Y$, define
\begin{equation}
  H(Y)={}_{\alpha}(eY).
  \label{eq:selection-functor}
\end{equation}
Here $eY$ is an $R$-submodule, and the subscript denotes restriction of
scalars along $\alpha$: on the vector space $eY$, an element $r\in R$
acts by the original operator $\alpha(r)$.  This construction defines an
additive functor on left $R$-modules and preserves finite dimensionality.
We write $H^{\circ j}$ for its $j$-fold iterate, with
$H^{\circ 0}$ the identity functor; superscripts without a circle on
$H$ will denote cohomology of a complex.

\begin{proposition}
\label{prop:selection}
There exist a finitely presented unital complex algebra $R$, a central
idempotent $e\in R$, and an automorphism $\alpha$ of $R$ such that, for
every integer $m\geq 1$, a finite-dimensional left $R$-module $Y_m$
satisfies
\begin{equation}
  H^{\circ j}(Y_m)\neq 0\quad(0\leq j<m),
  \qquad H^{\circ m}(Y_m)=0.
  \label{eq:extinction}
\end{equation}
The algebra, idempotent, and automorphism are independent of $m$.
\end{proposition}

We shall take $R=\C G$, where $G$ is a finitely presented group with
central involutions $z_N$ indexed by $N\in\Z$. An automorphism shifting
this family will give $\alpha$, and we shall take $e=(1+z_0)/2$.

The auxiliary group used below is motivated by finitely presented
matrix groups of Abels type. In particular, the commutator transfer
and centrality argument below parallels the proof of Proposition~4.9,
equations~(4.9)--(4.10), in \cite{Rego2022}. Our presentation and the
finite quotients used for selection are specified here, and all required
properties are proved directly.

\subsection{The group and its finite presentation}

Write $[x,y]=xyx^{-1}y^{-1}$ and let $I=\{1,2,3\}$.  Define a group
$G$ with commuting generators $T_\ell$, $\ell\in I$, and generators
\[
  U_i(r),\qquad V_i(r),\qquad W_{ij}(r)
  \quad(i,j\in I,\ i\neq j,\ r\in\Z).
\]
Each generator type $S$ has a weight $\lambda_S\in\Z^3$, given by
\begin{equation}
  \lambda_{U_i}=-\varepsilon_i,\qquad
  \lambda_{V_i}=\varepsilon_i,\qquad
  \lambda_{W_{ij}}=\varepsilon_i-\varepsilon_j,
  \label{eq:group-weights}
\end{equation}
where $\varepsilon_1,\varepsilon_2,\varepsilon_3$ is the standard basis.
Impose the conjugation rules
\begin{equation}
  T_\ell S(r)T_\ell^{-1}=S(r+\lambda_{S,\ell}).
  \label{eq:group-conjugation}
\end{equation}
The remaining relations, with all arguments ranging independently
over $\Z$, are
\begin{enumerate}
  \item $U_i(r)^2=1$;
  \item generators of types $U_i,U_j$, of types $V_i,V_j$, and of
    types $U_i,V_j$ commute whenever $i\neq j$;
  \item $W_{ij}$ commutes with $U_h$ for $h\neq i$ and with $V_h$
    for $h\neq j$;
  \item for $i\neq j$,
  \begin{equation}
    [U_i(r),W_{ij}(s)]=U_j(r+s),\qquad
    [W_{ij}(s),V_j(r)]=V_i(s+r).
    \label{eq:group-commutator-relations}
  \end{equation}
\end{enumerate}

\begin{lemma}
\label{lem:group-finite-presentation}
The group $G$ is finitely presented.
\end{lemma}

\begin{proof}
For $n=(n_1,n_2,n_3)\in\Z^3$, put
$T^n=T_1^{n_1}T_2^{n_2}T_3^{n_3}$ and regard a weight as the map
$\lambda_S(n)=\sum_\ell\lambda_{S,\ell}n_\ell$.  Retain only the
fifteen generators
\[
  T_\ell,\quad u_i=U_i(0),\quad v_i=V_i(0),\quad
  w_{ij}=W_{ij}(0)\quad(i\neq j).
\]
Besides commutation of the $T_\ell$, require the following finite
set of centralizer relations.  Each of $u_i,v_i$ commutes with
$T_h$ for $h\neq i$, and $w_{ij}$ commutes with $T_iT_j$ and with
$T_h$ for the unique $h\notin\{i,j\}$.  These are precisely
centralization by a basis of each kernel lattice:
\[
\begin{aligned}
 \ker\lambda_{U_i}=\ker\lambda_{V_i}
   &=\bigoplus_{h\neq i}\Z\varepsilon_h,\\
 \ker\lambda_{W_{ij}}
   &=\Z(\varepsilon_i+\varepsilon_j)
     \oplus\Z\varepsilon_h
       \quad(h\notin\{i,j\}).
\end{aligned}
\]
Since the $T_\ell$ commute, the kernel lattice also centralizes
every $T$-conjugate of its argument-zero generator.  Thus one may define
\begin{equation}
\begin{aligned}
  U_i(r)&=T_i^{-r}u_iT_i^r,&
  V_i(r)&=T_i^rv_iT_i^{-r},&
  W_{ij}(r)&=T_i^rw_{ij}T_i^{-r}.
\end{aligned}
\label{eq:group-finite-generators}
\end{equation}
Equivalently, $S(r)=T^nS(0)T^{-n}$ for any $n$ with
$\lambda_S(n)=r$.  The kernel relations make this independent of
the choice of $n$ and imply all of \eqref{eq:group-conjugation}.

It remains to impose the argument-zero instance of each of the
remaining relation patterns.  For every pair of input types actually
occurring in those patterns, the two weights give a surjection
$\Z^3\to\Z^2$.  The following table supplies integral preimages of
an arbitrary pair $(r,s)$; unmentioned coordinates are zero.
\begin{center}
\small
\begin{tabular}{@{}lll@{}}
\toprule
Input types & Index condition & Coordinates of $n$ \\
\midrule
$U_i,U_j$ & $i\neq j$ & $n_i=-r,\ n_j=-s$ \\
$V_i,V_j$ & $i\neq j$ & $n_i=r,\ n_j=s$ \\
$U_i,V_j$ & $i\neq j$ & $n_i=-r,\ n_j=s$ \\
$W_{ij},U_j$ & $i\neq j$ & $n_j=-s,\ n_i=r-s$ \\
$W_{ij},U_h$ & $h\notin\{i,j\}$ &
  $n_i=r,\ n_j=0,\ n_h=-s$ \\
$W_{ij},V_i$ & $i\neq j$ & $n_i=s,\ n_j=s-r$ \\
$W_{ij},V_h$ & $h\notin\{i,j\}$ &
  $n_i=r,\ n_j=0,\ n_h=s$ \\
$U_i,W_{ij}$ & $i\neq j$ & $n_i=-r,\ n_j=-r-s$ \\
$W_{ij},V_j$ & $i\neq j$ & $n_j=s,\ n_i=r+s$ \\
\bottomrule
\end{tabular}
\end{center}
Conjugating a zero-argument commutation relation by the corresponding
$T^n$ therefore gives every required instance.  For the two relations
in \eqref{eq:group-commutator-relations}, the output weight is the sum
of the input weights:
\[
  -\varepsilon_i+(\varepsilon_i-\varepsilon_j)=-\varepsilon_j,
  \qquad
  (\varepsilon_i-\varepsilon_j)+\varepsilon_j=\varepsilon_i.
\]
Hence the same conjugation also gives the correct output argument
$r+s$.  The square relation propagates from argument zero because
each individual weight is surjective onto $\Z$.

We have obtained all the defining relations from finitely many
relators on fifteen generators.  Conversely, the chosen kernel
centralizers and zero-argument relators hold in the original
presentation by \eqref{eq:group-conjugation}.  The two presentations
are therefore equivalent.
\end{proof}

\subsection{Central involutions and an automorphism}

The finite presentation is now fixed. We next identify central
involutions that the automorphism will shift; their independent images
in finite quotients will let us prescribe the extinction time.

\begin{lemma}
\label{lem:central-shift}
For each $N\in\Z$, the expression
\begin{equation}
  z_N=[U_i(a),V_i(b)],\qquad a+b=N,
  \label{eq:central-involutions}
\end{equation}
is independent of $i,a,b$.  Each $z_N$ is central and satisfies
$z_N^2=1$.  There is an automorphism $\alpha_G$ of $G$ fixing all
$T_\ell,V_i(r),W_{ij}(r)$ and sending $U_i(r)$ to $U_i(r+1)$;
it satisfies $\alpha_G(z_N)=z_{N+1}$.
\end{lemma}

\begin{proof}
Fix distinct $i,j$ and integers $r,q,s$, and set
\[
 x=U_i(r),\quad y=W_{ij}(q),\quad v=V_j(s),\qquad
 u=[x,y]=U_j(r+q),\quad w=[y,v]=V_i(q+s).
\]
The defining commutations give $[x,v]=[w,u]=[w,v]=1$.
Using $xy=uyx$ and $yv=wvy$, reorder $xyv$ in the two ways
\[
  xyv=uwvyx=(xwx^{-1})vuyx.
\]
With $c=[x,w]$, cancellation gives $uwv=cwvu$, and consequently
\[
  c=uwvu^{-1}v^{-1}w^{-1}=uvu^{-1}v^{-1}=[u,v].
\]
The last equality uses only the commutation of $w$ with $u$ and $v$.
We have proved
\begin{equation}
 [U_i(r),V_i(q+s)]=[U_j(r+q),V_j(s)].
 \label{eq:commutator-transfer}
\end{equation}
Given $a+b=c+d=N$ at distinct internal indices, choose
$r=a$, $s=d$, and $q=c-a=b-d$ in
\eqref{eq:commutator-transfer}.  This compares the two expressions
in \eqref{eq:central-involutions}.  Two splittings at the same
internal index can be compared through another index, proving
the claimed independence.

To show that $z_N$ commutes with $U_h(r)$ or $V_h(r)$, express
$z_N$ using an index $i\neq h$.  Both factors of that commutator
commute with the tested generator.  To commute with $W_{hj}(r)$,
use the index outside $\{h,j\}$, so that again both factors commute
with it.  Finally, conjugation by $T_\ell$ changes the arguments
$(a,b)$ in \eqref{eq:central-involutions} to
$(a-\delta_{\ell i},b+\delta_{\ell i})$.  Their sum is unchanged,
so the splitting independence proves invariance under $T_\ell$.
Thus $z_N$ is central.  Since $U_i(a)^2=1$, centrality also gives
\[
  1=[U_i(a)^2,V_i(b)]=[U_i(a),V_i(b)]^2=z_N^2.
\]

Translation of every $U$-argument by $1$ preserves the conjugation,
square, and commutation relations.  The first relation in
\eqref{eq:group-commutator-relations} becomes
\[
  [U_i(r+1),W_{ij}(s)]=U_j(r+s+1),
\]
which is another defining instance; the second relation is unchanged.
This defines an endomorphism $\alpha_G$.  Translation by $-1$
preserves the relations for the same reason and is its two-sided
inverse.  Formula \eqref{eq:central-involutions} then gives
$\alpha_G(z_N)=z_{N+1}$.
\end{proof}

\subsection{Finite quotients detecting independent involutions}

\begin{lemma}
\label{lem:finite-central-quotients}
For every $m\geq1$, there is a finite quotient $F_m$ of $G$ in which
the images of $z_0,\ldots,z_{m-1}$ generate a central subgroup
isomorphic to $(\Z/2\Z)^m$.
\end{lemma}

\begin{proof}
Let $S_m=\F_2[t]/(t^m-1)$.  This ring is finite, and $t$ is a unit
with inverse $t^{m-1}$, so $t^r$ is defined for every $r\in\Z$.
Use matrices indexed by $0,1,2,3,4$.  Send $T_\ell$ to the diagonal
matrix with $t$ in position $\ell$ and $1$ elsewhere, and set
\begin{equation}
\begin{aligned}
 U_i(r)&\longmapsto I_5+t^rE_{0i},&
 V_i(r)&\longmapsto I_5+t^rE_{i4},&
 W_{ij}(r)&\longmapsto I_5+t^rE_{ij}.
\end{aligned}
\label{eq:finite-matrix-images}
\end{equation}
Here $E_{ab}$ denotes a matrix unit.  Every displayed elementary
matrix is invertible, with itself as inverse: its off-diagonal
matrix unit has square zero, and the coefficient ring has
characteristic two.

We check the relations.  Diagonal conjugation multiplies $E_{ab}$
by the ratio of the diagonal entries at $a$ and $b$, which gives
exactly \eqref{eq:group-weights}.  The identity
$E_{ab}E_{cd}=\delta_{bc}E_{ad}$ shows that all products in both
orders vanish for each required commuting pair: two units $E_{0i}$,
two units $E_{i4}$, the pair $E_{0i},E_{j4}$ for $i\neq j$, the pair
$E_{ij},E_{0h}$ for $h\neq i$, and the pair $E_{ij},E_{h4}$ for
$h\neq j$.  The required square relations hold by the same
square-zero calculation.  For $i\neq j$, the nonzero product in the
first noncommutation relation is $E_{0i}E_{ij}=E_{0j}$, and in the
second it is $E_{ij}E_{j4}=E_{i4}$.  The higher products vanish,
giving
\[
\begin{aligned}
 [I_5+t^rE_{0i},I_5+t^sE_{ij}]&=I_5+t^{r+s}E_{0j},\\
 [I_5+t^sE_{ij},I_5+t^rE_{j4}]&=I_5+t^{s+r}E_{i4}.
\end{aligned}
\]
Thus \eqref{eq:finite-matrix-images} defines a homomorphism from
$G$ into the finite group $\operatorname{GL}_5(S_m)$.  Let $F_m$
be its image.

The image of $z_N$ is $I_5+t^NE_{04}$.  It is central in $F_m$:
the unit $E_{04}$ has zero product in both orders with every
off-diagonal unit in \eqref{eq:finite-matrix-images}, and the
diagonal matrices have entry $1$ at both $0$ and $4$.  For
$c_0,\ldots,c_{m-1}\in\{0,1\}$,
\[
 \prod_{q=0}^{m-1}(I_5+t^qE_{04})^{c_q}
   =I_5+\left(\sum_{q=0}^{m-1}c_qt^q\right)E_{04}.
\]
The classes $1,t,\ldots,t^{m-1}$ are an $\F_2$-basis of $S_m$,
because the defining polynomial is monic of degree $m$.
Thus this product is the identity only when all $c_q$ are zero.
This argument also applies when $m$ is even and $S_m$ is
nonreduced.  The claimed independence and orders follow.
\end{proof}

\subsection{The algebra and its test modules}

The quotient in \Cref{lem:finite-central-quotients} supplies the independent
signs. We pass to the group algebra and choose the signs so that each
selection before the last is the identity, while the last is zero.

\begin{proof}[Proof of \Cref{prop:selection}]
Take $R=\C G$.  A finite group presentation yields a finite
presentation of its unital group algebra: for each group generator
adjoin a formal inverse and the two inverse equations, and impose
the finitely many group relators as algebra relations.  Hence $R$
is finitely presented by \Cref{lem:group-finite-presentation}.
Extend $\alpha_G$ linearly to an algebra automorphism $\alpha$ and put
\[
  e=\frac{1+z_0}{2}.
\]
By \Cref{lem:central-shift}, this is a central idempotent.  For
$q\geq0$, write
\[
  e_q=\alpha^q(e)=\frac{1+z_q}{2}.
\]
These are commuting central idempotents of the one fixed algebra $R$.

For the finite quotient in \Cref{lem:finite-central-quotients}, let
$\overline z_q$ be the image of $z_q$ and let
$Z_m=\langle\overline z_0,\ldots,\overline z_{m-1}\rangle$.
Independence defines a character $\chi_m:Z_m\to\{1,-1\}$ by
\[
 \chi_m(\overline z_q)=
 \begin{cases}
   1,&0\leq q<m-1,\\
   -1,&q=m-1.
 \end{cases}
\]
In the complex regular representation of $F_m$, form
\[
  p_{\chi_m}=\frac{1}{|Z_m|}
       \sum_{z\in Z_m}\chi_m(z)^{-1}z,
  \qquad
  Y_m=\C[F_m]p_{\chi_m}.
\]
The element $p_{\chi_m}$ is nonzero, since its identity coefficient
is $1/|Z_m|$, and $z p_{\chi_m}=\chi_m(z)p_{\chi_m}$ for
$z\in Z_m$.  Centrality of $Z_m$ then shows that every vector of
$Y_m$ has this same character.  In particular, $Y_m$ is a nonzero
finite-dimensional left $\C[F_m]$-module.  Regard it as an
$R$-module through the quotient map.

For every left $R$-module $Y$, induction in the definition of $H$
gives an identification
\begin{equation}
  H^{\circ j}(Y)\cong {}_{\alpha^j}
       \bigl(e_0e_1\cdots e_{j-1}Y\bigr),
  \qquad j\geq0,
  \label{eq:selection-iterates}
\end{equation}
where the empty product is $1$.  Indeed, on the vector space at
stage $j$, selection by $e$ acts as the original operator
$\alpha^j(e)=e_j$.  Pulling the resulting action back once more
along $\alpha$ makes $r$ act by $\alpha^{j+1}(r)$, proving the
inductive step.  Centrality of all the $e_q$ makes the subspace in
\eqref{eq:selection-iterates} an $R$-submodule before this pullback.

On $Y_m$, the operators $e_0,\ldots,e_{m-2}$ are the identity and
$e_{m-1}$ is zero.  Consequently \eqref{eq:selection-iterates}
gives \eqref{eq:extinction}, also for $m=1$.  This argument does
not require $\alpha$ to preserve the kernel of the finite quotient:
each pullback uses the same finite-dimensional vector space with
its action precomposed by an automorphism of $R$.  The construction
of $R,e,\alpha$ preceded the choice of $m$, as required.
\end{proof}
